\chapter{Introduction}\label{ch:intro}

\section{Hallucination in Large Language Models}\label{sec:hallucination}
\subsection{Where Hallucination Originates: Training Data, Model, and Prompt}
\subsection{Why Multi-Hop Questions Are the Hard Case}

\section{From Retrieval Augmentation to Graph-Structured Grounding}\label{sec:rag-to-graph}
\subsection{Vector Retrieval and Its Structural Blind Spot}
\subsection{Knowledge Graphs and Static GraphRAG}
\subsection{Agentic Retrieval: Reasoning as Navigation}

\section{Limitations of Existing Agentic KGQA Systems}\label{sec:limitations}
% Frame 1.3.1 as a PRECISION deficit, not a groundedness one. Prior agentic
% systems emit whatever the final LLM call produces from the traversed context,
% with no check that each asserted entity is supported AS AN ANSWER. Point
% forward to the precision gap in Chapter~\ref{ch:results}. Do NOT claim prior
% systems hallucinate structurally --- Section~\ref{sec:groundedness-results}
% shows they do not.
\subsection{No Check on What the Final Answer Asserts}
\subsection{Unquantified Contribution of Individual Agentic Mechanisms}
\subsection{Accuracy Reported Without Cost}

\section{Problem Statement}\label{sec:problem}

\section{Research Questions}\label{sec:rqs}
\subsection{RQ1: Does Agentic Navigation Improve Multi-Hop Factual Accuracy?}
\subsection{RQ2: What Does Pre-Generation Verification Contribute Beyond Graph
  Navigation?}
\subsection{RQ3: Which Components Contribute What, at What Token Cost?}

\section{Our Contribution}\label{sec:contribution}
\subsection{The AGR Framework and Its Verification Layer}
\subsection{A Component-Level Ablation of Agentic Mechanisms}
\subsection{Stratum-Dependent Decomposition: When Planning Hurts}
\subsection{The Echo Attractor as a Named Failure Mode}
\subsection{Quantified Benchmark Defect Rates for WebQSP and CWQ}
\subsection{An Evaluation Protocol with Pre-Registered Thresholds}

\section{Scope and Delimitations}\label{sec:scope}

\section{Thesis Organization}\label{sec:organization}

\endinput
